The covariance estimates must now produce something random: an operator path
driven by the noise of the original equation. Another Gaussian field with a
convenient covariance would not do, because the later comparisons are made
on the original noise (Figure~\ref{Fig:common-noise}).

The wave phase is a displacement in $H^{1/2}$ and a velocity
in $H^{-1/2}$. For a fixed spatial mask and a finite time horizon,
Section~\ref{S:actual-operator-driver} builds the integrated interaction path
on this space. Proposition~\ref{P:actual-operator-driver} states its regularity
and how it is approximated. Section~\ref{S:source-driver-convergence} matches
the regularized sources to the original noise and proves convergence in the
path norm. The construction keeps the joint noise information that Gaussian
changes of law will need.

The chapter uses the two variance directions and the compact
Gaussian estimates of Chapter~\ref{S:book-gaussian-foundations}. On
a first reading, keep Propositions~\ref{P:actual-operator-driver}
and~\ref{P:source-gaussian-transfer} and let the proofs wait. The
dependence of the operator-path bounds on the mask radius is tracked in
Chapter~\ref{S:book-growing-radius}. Episode~\VideoRef{5}{FGR9ShF2F0w}
of the video series covers this chapter.
